\documentclass[11pt,a4paper]{article}

% --- Pacotes de Idioma e Codificação ---
\usepackage[T1]{fontenc}
\usepackage[portuguese]{babel}

% --- Design e Layout ---
\usepackage[margin=2.5cm]{geometry}
\usepackage{charter}
\usepackage{lastpage}
\usepackage{fancyhdr}
\usepackage[table]{xcolor}
\usepackage{tabularx}
\usepackage{booktabs}
\usepackage{enumitem}
\usepackage{amssymb}
\usepackage{parskip}

% --- Configuração de Cores ---
\definecolor{primaryblue}{RGB}{0, 51, 102}
\definecolor{sectionblue}{RGB}{0, 102, 153}

% --- Cabeçalho e Rodapé ---
\pagestyle{fancy}
\fancyhf{}
\fancyhead[L]{\small \textbf{Requisitos Reformatados} | EcoRide Solutions}
\fancyhead[R]{\small \thepage\ de \pageref{LastPage}}
\fancyfoot[L]{\small \textit{Confidencial - Apenas para uso interno}}
\renewcommand{\headrulewidth}{0.4pt}

% --- Estilos Customizados ---
\newcommand{\reqheader}[1]{{\noindent\large\bfseries\color{sectionblue}#1}\par\vspace{0.3cm}}
\newcommand{\reqfield}[2]{{\noindent\textbf{#1:} #2}\par\vspace{0.2cm}}
\newcommand{\reqsystem}[1]{
	\begin{itemize}[leftmargin=1cm]
		#1
	\end{itemize}
}

% --- Início do Documento ---
\begin{document}
	
	{\noindent\LARGE\bfseries\color{primaryblue} Requisitos do Sistema EcoRide Solutions}
	
	{\noindent\large\textit{Formato Estruturado de Requisitos}}\par\vspace{1cm}

% ======================== REQUISITOS GERAIS ========================

\section*{Requisitos Gerais}

% --- Requisito Geral 1 ---
\reqheader{REQ-GEN-001: Cálculo Automático do Valor Final da Reparação}

\reqfield{Requisito de utilizador}{O sistema é capaz de calcular automaticamente o valor final de uma reparação}

\reqfield{Fonte}{Documento de Análise de Requisitos}

\reqfield{Área do sistema}{Gestão de Ordens de Serviço, Gestão Financeira}

\reqfield{Requisitos de sistema}{
	\begin{enumerate}
		\item O sistema deve somar automaticamente o custo de todas as reparações realizadas
		\item O sistema deve somar automaticamente o custo de todas as peças utilizadas
		\item O preço unitário de cada reparação é definido previamente pelo gerente
		\item O preço unitário de cada peça é definido previamente na ficha da peça
		\item O valor final é apresentado na vista de "Pendente Pagamento" da OS
	\end{enumerate}
}

\reqfield{Relevância para a aplicação}{O cálculo automático garante precisão financeira, elimina erros manuais e agiliza o processo de cobrança aos clientes}

\vspace{0.8cm}

% --- Requisito Geral 2 ---
\reqheader{REQ-GEN-002: Taxa de Armazenagem Automática por Atraso}

\reqfield{Requisito de utilizador}{O sistema é capaz de aplicar automaticamente uma taxa de armazenagem quando um cliente não levanta a trotinete dentro do prazo}

\reqfield{Fonte}{Documento de Análise de Requisitos}

\reqfield{Área do sistema}{Gestão de Ordens de Serviço, Gestão Financeira}

\reqfield{Requisitos de sistema}{
	\begin{enumerate}
		\item O sistema registar a data de notificação do cliente sobre conclusão da reparação
		\item A partir do 7º dia após notificação, é aplicada uma taxa diária de 3€
		\item A taxa é acumulada por cada dia completo decorrido
		\item O valor máximo de taxa acumulada é 90€ (correspondendo a 30 dias)
		\item O valor total com taxa é apresentado na OS em estado "Pendente Pagamento"
	\end{enumerate}
}

\reqfield{Relevância para a aplicação}{Incentiva o cliente a levantar a trotinete rapidamente, evitando custo de armazenagem desnecessário para a oficina}

\vspace{0.8cm}

% --- Requisito Geral 3 ---
\reqheader{REQ-GEN-003: Estimativa de Tempo para Conclusão da Reparação}

\reqfield{Requisito de utilizador}{O sistema é capaz de estimar automaticamente quanto tempo levará a reparação de uma trotinete}

\reqfield{Fonte}{Documento de Análise de Requisitos}

\reqfield{Área do sistema}{Gestão de Ordens de Serviço, Tabela de Reparações}

\reqfield{Requisitos de sistema}{
	\begin{enumerate}
		\item O sistema consulta o tempo médio registado na tabela de reparações para cada tipo de reparação
		\item O sistema contabiliza todas as reparações necessárias da OS atual
		\item O sistema contabiliza o número total de OS pendentes de diagnóstico e reparação
		\item A estimativa é calculada como: (tempo total das reparações da OS) + (tempo de espera das OS pendentes)
		\item A estimativa é apresentada no momento de criação da OS
	\end{enumerate}
}

\reqfield{Relevância para a aplicação}{Permite ao cliente ter expectativa realista sobre quando sua trotinete estará pronta, melhorando a experiência e comunicação}

\vspace{0.8cm}

% --- Requisito Geral 4 ---
\reqheader{REQ-GEN-004: Rastreabilidade de Associação de Peças}

\reqfield{Requisito de utilizador}{O sistema regista quem associou uma peça a uma OS e quando foi feito}

\reqfield{Fonte}{Documento de Análise de Requisitos}

\reqfield{Área do sistema}{Gestão de Ordens de Serviço, Auditoria}

\reqfield{Requisitos de sistema}{
	\begin{enumerate}
		\item Quando uma peça é associada a uma OS, o sistema registar automaticamente o identificador do utilizador
		\item O sistema registar automaticamente a data e hora exata da associação
		\item Este registo é imutável e permanece no histórico da OS
		\item O histórico de associações é visível na ficha detalhada da OS
	\end{enumerate}
}

\reqfield{Relevância para a aplicação}{Garante rastreabilidade completa para auditoria, responsabilização e resolução de disputas sobre quem realizou ação}

\vspace{0.8cm}

% --- Requisito Geral 5 ---
\reqheader{REQ-GEN-005: Ordenação de OS por Estado Respeitando Data de Criação}

\reqfield{Requisito de utilizador}{O sistema organiza as ordens de serviço em listas por estado, mantendo ordem cronológica}

\reqfield{Fonte}{Documento de Análise de Requisitos}

\reqfield{Área do sistema}{Gestão de Ordens de Serviço}

\reqfield{Requisitos de sistema}{
	\begin{enumerate}
		\item Cada estado da OS tem uma lista própria (Pendente Diagnóstico, Pendente Reparação, etc.)
		\item Quando uma OS muda de estado, é movida para a lista do novo estado
		\item Dentro de cada lista, as OS são ordenadas pela data de criação (mais antigas primeiro)
		\item A ordenação é mantida dinamicamente conforme novos estados são atingidos
	\end{enumerate}
}

\reqfield{Relevância para a aplicação}{Garante que as reparações mais antigas sejam atendidas primeiro (FIFO), evitando que clientes esperem indefinidamente}

\vspace{0.8cm}

% --- Requisito Geral 6 ---
\reqheader{REQ-GEN-006: Alertas Automáticos para Secretária}

\reqfield{Requisito de utilizador}{O sistema notifica automaticamente a secretária quando uma OS atinge certos estados críticos}

\reqfield{Fonte}{Documento de Análise de Requisitos}

\reqfield{Área do sistema}{Gestão de Ordens de Serviço, Gestão de Alertas}

\reqfield{Requisitos de sistema}{
	\begin{enumerate}
		\item Quando uma OS é colocada no estado "Pendente de aprovação de orçamento", um alerta é gerado para a secretária
		\item Quando uma OS é colocada no estado "Pendente Pagamento", um alerta é gerado para a secretária
		\item Quando uma OS é colocada no estado "A aguardar peças", um alerta é gerado para a secretária
		\item Cada alerta contém a referência única da OS e o estado atual
		\item Cada alerta tem um campo que indica se foi "tratado" ou não
		\item Os alertas são armazenados na lista pessoal de alertas da secretária
		\item A secretária pode marcar manualmente um alerta como tratado
	\end{enumerate}
}

\reqfield{Relevância para a aplicação}{Garante que ações críticas não sejam esquecidas, reduzindo atrasos e melhorando o fluxo de trabalho}

\vspace{0.8cm}

% --- Requisito Geral 7 ---
\reqheader{REQ-GEN-007: Rastreabilidade de Responsáveis por Passo da Execução}

\reqfield{Requisito de utilizador}{O sistema regista qual mecânico realizou cada fase da reparação}

\reqfield{Fonte}{Documento de Análise de Requisitos}

\reqfield{Área do sistema}{Gestão de Ordens de Serviço, Auditoria}

\reqfield{Requisitos de sistema}{
	\begin{enumerate}
		\item Quando um mecânico completa o diagnóstico, o sistema registar quem realizou
		\item Quando um mecânico conclui a reparação, o sistema registar quem realizou
		\item Se múltiplas reparações são realizadas, cada uma registar o mecânico responsável
		\item Estes registos são permanentes e visíveis no histórico da OS
	\end{enumerate}
}

\reqfield{Relevância para a aplicação}{Permite responsabilização dos mecânicos e rastreamento de qualidade de trabalho para fins de avaliação e melhoria}

\vspace{0.8cm}

% --- Requisito Geral 8 ---
\reqheader{REQ-GEN-008: Bloqueio de Atribuição Simultânea de OS}

\reqfield{Requisito de utilizador}{O sistema impede que dois mecânicos trabalhem simultaneamente na mesma ordem de serviço}

\reqfield{Fonte}{Documento de Análise de Requisitos}

\reqfield{Área do sistema}{Gestão de Ordens de Serviço, Controlo de Acesso}

\reqfield{Requisitos de sistema}{
	\begin{enumerate}
		\item Quando um mecânico seleciona uma OS para trabalhar, a OS é "bloqueada" para outros mecânicos
		\item Se outro mecânico tentar selecionar a mesma OS, recebe uma mensagem de erro
		\item A OS mantém-se bloqueada até o primeiro mecânico a libertar (marcar como concluída ou parada)
		\item O sistema mostra qual mecânico tem a OS atualmente bloqueada
	\end{enumerate}
}

\reqfield{Relevância para a aplicação}{Evita conflitos de trabalho duplicado, corrupção de dados e confusão sobre quem está responsável pela reparação}

\vspace{0.8cm}

% --- Requisito Geral 9 ---
\reqheader{REQ-GEN-009: Gestão de Stock Automática por Peças Associadas}

\reqfield{Requisito de utilizador}{O sistema atualiza automaticamente o stock quando peças são associadas ou removidas de uma OS}

\reqfield{Fonte}{Documento de Análise de Requisitos}

\reqfield{Área do sistema}{Gestão de Stock, Gestão de Ordens de Serviço}

\reqfield{Requisitos de sistema}{
	\begin{enumerate}
		\item Quando um mecânico associa uma peça a uma OS em "Pendente Reparação", a quantidade em stock diminui 1
		\item Quando uma secretária associa uma peça a uma OS em "Pendente Pagamento", a quantidade em stock diminui 1
		\item Quando um mecânico revoga uma peça da OS, a quantidade em stock aumenta 1
		\item O stock é atualizado imediatamente, evitando overbooking
		\item A quantidade em stock nunca pode ficar negativa
	\end{enumerate}
}

\reqfield{Relevância para a aplicação}{Garante que o stock sempre reflete a realidade, evitando reservas falsas e facilitando encomendas adequadas}

\vspace{0.8cm}

% --- Requisito Geral 10 ---
\reqheader{REQ-GEN-010: Alertas Automáticos para Gestor de Stock}

\reqfield{Requisito de utilizador}{O sistema notifica automaticamente o gestor de stock quando há problemas com peças}

\reqfield{Fonte}{Documento de Análise de Requisitos}

\reqfield{Área do sistema}{Gestão de Stock, Gestão de Alertas}

\reqfield{Requisitos de sistema}{
	\begin{enumerate}
		\item Quando uma peça é marcada como "Possível defeito", um alerta é gerado para o gestor de stock
		\item Quando o stock de uma peça atinge o nível mínimo, um alerta é gerado para o gestor de stock
		\item Cada alerta contém o motivo (defeito ou stock baixo) e a peça em questão
		\item Cada alerta tem um campo que indica se foi "tratado" ou não
		\item Os alertas são armazenados na lista pessoal de alertas do gestor de stock
		\item O gestor de stock pode marcar manualmente um alerta como tratado
	\end{enumerate}
}

\reqfield{Relevância para a aplicação}{Reduz tempo de reação a problemas de stock, evitando paragens por falta de peças e detectando peças defeituosas}

\vspace{0.8cm}

% --- Requisito Geral 11 ---
\reqheader{REQ-GEN-011: Auditoria Completa de Movimentos de Stock}

\reqfield{Requisito de utilizador}{O sistema mantém um registro completo de todas as operações de stock}

\reqfield{Fonte}{Documento de Análise de Requisitos}

\reqfield{Área do sistema}{Gestão de Stock, Auditoria}

\reqfield{Requisitos de sistema}{
	\begin{enumerate}
		\item Cada entrada de peças é registada com data, hora e utilizador responsável
		\item Cada saída de peças é registada com data, hora e utilizador responsável
		\item Alterações aos dados de uma peça (preço, fornecedor, etc.) são registadas com data, hora e utilizador
		\item Alterações aos dados de um fornecedor são registadas com data, hora e utilizador
		\item Este log de auditoria é imutável e permanente
		\item O log é acessível apenas ao gestor de stock e gerente
	\end{enumerate}
}

\reqfield{Relevância para a aplicação}{Fornece rastreabilidade completa para auditoria financeira, conformidade regulatória e investigação de discrepâncias}

\vspace{0.8cm}

% --- Requisito Geral 12 ---
\reqheader{REQ-GEN-012: Interface Personalizada para Mecânico}

\reqfield{Requisito de utilizador}{A interface do mecânico apresenta apenas informação relevante ao seu trabalho}

\reqfield{Fonte}{Documento de Análise de Requisitos}

\reqfield{Área do sistema}{Interface de Utilizador, Gestão de Ordens de Serviço}

\reqfield{Requisitos de sistema}{
	\begin{enumerate}
		\item A interface mostra a informação pessoal do mecânico (nome, identificador)
		\item A interface mostra uma lista de OS em "Pendente Diagnóstico" disponíveis
		\item A interface mostra uma lista de OS em "Pendente Reparação" disponíveis
		\item Quando um mecânico está a trabalhar em uma OS, mostra os detalhes completos dessa OS
		\item A interface mostra apenas as ações e opções aplicáveis ao estado atual da OS
		\item Nenhuma informação de outros utilizadores ou OS é exibida por padrão
	\end{enumerate}
}

\reqfield{Relevância para a aplicação}{Melhora a usabilidade e foco do mecânico, evitando confusão e reduzindo tempo de procura de informação}

\vspace{0.8cm}

% --- Requisito Geral 13 ---
\reqheader{REQ-GEN-013: Saída de Stock para Peça em Devolução}

\reqfield{Requisito de utilizador}{O stock de peças é atualizado automaticamente quando uma peça é marcada como possível defeito e entra em processo de devolução}

\reqfield{Fonte}{Documento de Análise de Requisitos}

\reqfield{Área do sistema}{Gestão de Stock, Gestão de Peças Defeituosas}

\reqfield{Requisitos de sistema}{
	\begin{enumerate}
		\item Quando uma peça "Possível defeito" é colocada em estado "Pendente devolução", a quantidade em stock diminui 1
		\item A peça deixa de estar disponível para uso em novas OS
		\item O movimento é registado na auditoria de stock
	\end{enumerate}
}

\reqfield{Relevância para a aplicação}{Mantém o stock consistente ao remover peças defeituosas de circulação}

\vspace{0.8cm}

% --- Requisito Geral 14 ---
\reqheader{REQ-GEN-014: Entrada de Stock para Peça Devolvida}

\reqfield{Requisito de utilizador}{O stock de peças é aumentado automaticamente quando uma devolução ao fornecedor é concluída com sucesso}

\reqfield{Fonte}{Documento de Análise de Requisitos}

\reqfield{Área do sistema}{Gestão de Stock, Devolução a Fornecedores}

\reqfield{Requisitos de sistema}{
	\begin{enumerate}
		\item Quando uma peça em estado "Pendente devolução" é colocada em estado "Devolvida ao fornecedor", a quantidade em stock aumenta 1
		\item A peça volta a estar disponível para uso (se for recondicionada) ou é completamente removida conforme política
		\item O movimento é registado na auditoria de stock
	\end{enumerate}
}

\reqfield{Relevância para a aplicação}{Recupera valor do stock quando fornecedor aceita devolução, compensando custo de peças defeituosas}

\vspace{0.8cm}

% --- Requisito Geral 15 ---
\reqheader{REQ-GEN-015: Registro de Custo para Peça Inválida para Devolução}

\reqfield{Requisito de utilizador}{O sistema registar automaticamente o custo financeiro quando uma peça defeituosa não pode ser devolvida}

\reqfield{Fonte}{Documento de Análise de Requisitos}

\reqfield{Área do sistema}{Gestão de Stock, Gestão Financeira}

\reqfield{Requisitos de sistema}{
	\begin{enumerate}
		\item Quando uma peça em estado "Pendente devolução" é colocada em estado "Inválida para devolução", o custo é registado
		\item O custo adicionado é o valor de compra original da peça
		\item Este custo é adicionado ao "Gasto com peças" do mês em questão
		\item O movimento aparece no relatório financeiro do mês
	\end{enumerate}
}

\reqfield{Relevância para a aplicação}{Garante que perdas por peças defeituosas não retornáveis sejam refletidas nos gastos mensais para análise financeira}

\vspace{1.5cm}

% ======================== REQUISITOS DO GERENTE ========================

\section*{Requisitos do Gerente}

% --- Requisito Gerente 1 ---
\reqheader{REQ-GER-001: Acesso Administrativo Total}

\reqfield{Requisito de utilizador}{O gerente tem acesso a todas as funcionalidades do sistema}

\reqfield{Fonte}{Documento de Análise de Requisitos}

\reqfield{Área do sistema}{Controlo de Acesso}

\reqfield{Requisitos de sistema}{
	\begin{enumerate}
		\item O gerente pode executar todas as ações disponíveis para a secretária
		\item O gerente pode executar todas as ações disponíveis para o gestor de stock
		\item O gerente pode executar todas as ações disponíveis para os mecânicos
		\item O gerente tem permissões adicionais exclusivas (gestão de funcionários, preços, etc.)
	\end{enumerate}
}

\reqfield{Relevância para a aplicação}{Garante que o gerente pode intervir em qualquer situação operacional e tem visibilidade total da empresa}

\vspace{0.8cm}

% --- Requisito Gerente 2 ---
\reqheader{REQ-GER-002: Gestão de Funcionários}

\reqfield{Requisito de utilizador}{O gerente é capaz de adicionar, remover e modificar registos de funcionários}

\reqfield{Fonte}{Documento de Análise de Requisitos}

\reqfield{Área do sistema}{Gestão de Recursos Humanos}

\reqfield{Requisitos de sistema}{
	\begin{enumerate}
		\item O gerente pode criar novo registo de funcionário
		\item O gerente pode editar dados existentes de um funcionário
		\item O gerente pode desativar ou remover um funcionário
		\item O sistema mantém histórico de alterações (quando foi demitido, etc.)
	\end{enumerate}
}

\reqfield{Relevância para a aplicação}{Centraliza a gestão de recursos humanos no sistema, evitando registos dispersos e facilitando auditoria}

\vspace{0.8cm}

% --- Requisito Gerente 3 ---
\reqheader{REQ-GER-003: Estrutura Completa de Dados de Funcionário}

\reqfield{Requisito de utilizador}{O sistema armazena toda informação pessoal e profissional de um funcionário}

\reqfield{Fonte}{Documento de Análise de Requisitos}

\reqfield{Área do sistema}{Gestão de Recursos Humanos}

\reqfield{Requisitos de sistema}{
	\begin{enumerate}
		\item Campo obrigatório: Nome completo
		\item Campo obrigatório: Morada
		\item Campo obrigatório: Telemóvel
		\item Campo obrigatório: Email
		\item Campo obrigatório: Data de nascimento
		\item Campo obrigatório: NISS (Número de Identificação de Segurança Social)
		\item Campo obrigatório: NIF (Número de Identificação Fiscal)
		\item Campo obrigatório: NUS (Número Único de Segurança)
		\item Campo obrigatório: IBAN (para transferências salariais)
		\item Campo obrigatório: Salário base mensal
		\item Campo opcional: Foto de identificação
		\item Campo opcional: Data de admissão
	\end{enumerate}
}

\reqfield{Relevância para a aplicação}{Garante que toda informação necessária para folha de pagamento, segurança social e conformidade legal está centralizada}

\vspace{0.8cm}

% --- Requisito Gerente 4 ---
\reqheader{REQ-GER-004: Registro e Cálculo de Horas Extraordinárias}

\reqfield{Requisito de utilizador}{O sistema permite registar horas extraordinárias e calcula automaticamente o valor a pagar}

\reqfield{Fonte}{Documento de Análise de Requisitos}

\reqfield{Área do sistema}{Gestão de Recursos Humanos, Folha de Pagamento}

\reqfield{Requisitos de sistema}{
	\begin{enumerate}
		\item O gerente pode registar horas extraordinárias por funcionário e por mês
		\item O sistema calcula automaticamente o custo das horas extraordinárias
		\item O cálculo é baseado no valor/hora (salário base / 160 horas típicas)
		\item Suplemento por hora extraordinária é aplicado (ex: 1.25x)
		\item O valor total a pagar no mês = salário base + custo horas extraordinárias
		\item Este cálculo é refletido no relatório de folha de pagamento
	\end{enumerate}
}

\reqfield{Relevância para a aplicação}{Automatiza cálculo trabalhoso, reduz erros e garante cumprimento de obrigações legais de remuneração}

\vspace{0.8cm}

% --- Requisito Gerente 5 ---
\reqheader{REQ-GER-005: Consulta de Estado Financeiro da Empresa}

\reqfield{Requisito de utilizador}{O gerente é capaz de visualizar todos os movimentos financeiros realizados}

\reqfield{Fonte}{Documento de Análise de Requisitos}

\reqfield{Área do sistema}{Gestão Financeira}

\reqfield{Requisitos de sistema}{
	\begin{enumerate}
		\item O sistema apresenta uma página com histórico completo de movimentos
		\item Os movimentos incluem: pagamentos de encomendas a fornecedores (débito)
		\item Os movimentos incluem: pagamento de salários (débito)
		\item Os movimentos incluem: receção de pagamentos de clientes (crédito)
		\item Cada movimento mostra: data, tipo, descrição, valor e saldo atual
		\item Os movimentos são ordenados por data (mais recentes primeiro)
	\end{enumerate}
}

\reqfield{Relevância para a aplicação}{Fornece visibilidade financeira real à gestão, essencial para decisões estratégicas e conformidade}

\vspace{0.8cm}

% --- Requisito Gerente 6 ---
\reqheader{REQ-GER-006: Filtragem e Análise de Movimentos Financeiros}

\reqfield{Requisito de utilizador}{O gerente pode filtrar movimentos financeiros para análise segmentada}

\reqfield{Fonte}{Documento de Análise de Requisitos}

\reqfield{Área do sistema}{Gestão Financeira}

\reqfield{Requisitos de sistema}{
	\begin{enumerate}
		\item O sistema permite filtrar por intervalo de datas (data inicial e final)
		\item O sistema permite filtrar por tipo: salários, gastos com peças, lucro mão de obra, lucro venda peças
		\item O sistema calcula e apresenta totalizadores por categoria filtrada
		\item O sistema permite múltiplos filtros simultâneos
		\item Os resultados podem ser exportados para análise externa
	\end{enumerate}
}

\reqfield{Relevância para a aplicação}{Permite análise profunda do negócio, identificar tendências, rentabilidade por área e tomar decisões informadas}

\vspace{0.8cm}

% --- Requisito Gerente 7 ---
\reqheader{REQ-GER-007: Gestão da Tabela de Reparações}

\reqfield{Requisito de utilizador}{O gerente pode adicionar, editar e remover tipos de reparação e seus preços}

\reqfield{Fonte}{Documento de Análise de Requisitos}

\reqfield{Área do sistema}{Tabela de Reparações}

\reqfield{Requisitos de sistema}{
	\begin{enumerate}
		\item O gerente pode criar novo tipo de reparação
		\item O gerente pode editar dados existentes de uma reparação
		\item O gerente pode remover ou desativar um tipo de reparação
		\item O sistema impede remoção se reparação foi usada em OS (apenas desativa)
		\item Alterações aparecem imediatamente para criação de novas OS
	\end{enumerate}
}

\reqfield{Relevância para a aplicação}{Mantém a tabela de preços atualizada conforme custos e mercado mudam, refletindo margem de lucro desejada}

\vspace{0.8cm}

% --- Requisito Gerente 8 ---
\reqheader{REQ-GER-008: Estrutura Completa de Entrada de Reparação}

\reqfield{Requisito de utilizador}{O sistema armazena toda informação necessária sobre um tipo de reparação}

\reqfield{Fonte}{Documento de Análise de Requisitos}

\reqfield{Área do sistema}{Tabela de Reparações}

\reqfield{Requisitos de sistema}{
	\begin{enumerate}
		\item Campo obrigatório: Referência única (ex: REP-001)
		\item Campo obrigatório: Nomenclatura (nome curto da reparação)
		\item Campo obrigatório: Descrição detalhada (instruções de execução)
		\item Campo obrigatório: Preço de execução (em euros)
		\item Campo opcional: Tempo médio estimado
		\item Campo opcional: Ferramentas necessárias
		\item Campo opcional: Normas de segurança aplicáveis
	\end{enumerate}
}

\reqfield{Relevância para a aplicação}{Padroniza a execução de reparações e garante que prços cobrados cobrem tempo e recursos investidos}

\vspace{1.5cm}

% ======================== REQUISITOS DO GESTOR DE STOCK ========================

\section*{Requisitos do Gestor de Stock}

% --- Requisito Stock 1 ---
\reqheader{REQ-STK-001: Consulta Multi-Estados de Peças}

\reqfield{Requisito de utilizador}{O gestor de stock pode visualizar peças em diferentes estados do ciclo de vida}

\reqfield{Fonte}{Documento de Análise de Requisitos}

\reqfield{Área do sistema}{Gestão de Stock}

\reqfield{Requisitos de sistema}{
	\begin{enumerate}
		\item O sistema mostra lista de peças que existem no catálogo (nunca compradas)
		\item O sistema mostra lista de peças em stock (disponíveis)
		\item O sistema mostra lista de peças com possível defeito
		\item O sistema mostra lista de peças pendentes de devolução (saídas recentes)
		\item O sistema mostra lista de peças devolvidas ao fornecedor (devolução concluída)
		\item Cada lista é visível em separadores ou abas distintas
	\end{enumerate}
}

\reqfield{Relevância para a aplicação}{Fornece visibilidade do status de cada peça, facilitando gestão de qualidade e fluxo de fornecimento}

\vspace{0.8cm}

% --- Requisito Stock 2 ---
\reqheader{REQ-STK-002: Filtragem Avançada de Peças}

\reqfield{Requisito de utilizador}{O gestor de stock pode filtrar peças por múltiplos critérios}

\reqfield{Fonte}{Documento de Análise de Requisitos}

\reqfield{Área do sistema}{Gestão de Stock}

\reqfield{Requisitos de sistema}{
	\begin{enumerate}
		\item Filtro por estado (catálogo, stock, defeito, devolução, devolvido)
		\item Filtro por fornecedor (nome)
		\item Filtro por referência (número do fornecedor ou interna)
		\item Filtro por data de receção (peças recebidas entre datas)
		\item Filtro por tipo/categoria da peça
		\item Filtro por preço (peças acima/abaixo de valor)
		\item Combinação de múltiplos filtros
	\end{enumerate}
}

\reqfield{Relevância para a aplicação}{Agiliza a procura de peças específicas, economizando tempo do gestor de stock}

\vspace{0.8cm}

% --- Requisito Stock 3 ---
\reqheader{REQ-STK-003: Gestão de Quantidade em Stock}

\reqfield{Requisito de utilizador}{O gestor de stock pode adicionar, remover e editar quantidades de peças}

\reqfield{Fonte}{Documento de Análise de Requisitos}

\reqfield{Área do sistema}{Gestão de Stock}

\reqfield{Requisitos de sistema}{
	\begin{enumerate}
		\item O gestor pode ajustar manualmente a quantidade de uma peça em stock
		\item O gestor pode fazer contagem de stock e corrigir discrepâncias
		\item O sistema mantém histórico de ajustes manuais (data, utilizador, motivo)
		\item Alertas são disparados se quantidade fica abaixo do nível mínimo
	\end{enumerate}
}

\reqfield{Relevância para a aplicação}{Permite correção de erros de contagem e ajustes de stock por perdas ou danos}

\vspace{0.8cm}

% --- Requisito Stock 4 ---
\reqheader{REQ-STK-004: Gestão de Fornecedores}

\reqfield{Requisito de utilizador}{O gestor de stock pode adicionar, editar e remover fornecedores}

\reqfield{Fonte}{Documento de Análise de Requisitos}

\reqfield{Área do sistema}{Gestão de Fornecedores}

\reqfield{Requisitos de sistema}{
	\begin{enumerate}
		\item O gestor pode criar novo registo de fornecedor
		\item O gestor pode editar dados de fornecedor existente
		\item O gestor pode desativar um fornecedor
		\item O sistema impede remoção de fornecedor se tem peças ativas (apenas desativa)
		\item Histórico de alterações é mantido
	\end{enumerate}
}

\reqfield{Relevância para a aplicação}{Centraliza dados de fornecedores para facilitar comunicação, encomendas e rastreamento de cumprimento}

\vspace{0.8cm}

% --- Requisito Stock 5 ---
\reqheader{REQ-STK-005: Estrutura de Dados de Fornecedor}

\reqfield{Requisito de utilizador}{O sistema armazena informação essencial de cada fornecedor}

\reqfield{Fonte}{Documento de Análise de Requisitos}

\reqfield{Área do sistema}{Gestão de Fornecedores}

\reqfield{Requisitos de sistema}{
	\begin{enumerate}
		\item Campo obrigatório: Nome da empresa fornecedora
		\item Campo obrigatório: Email de contacto
		\item Campo obrigatório: Número de telefone de contacto
		\item Campo opcional: Prazo de entrega típico (em dias)
		\item Campo opcional: Morada
		\item Campo opcional: NIF/CNPJ
		\item Campo opcional: Condições de pagamento
	\end{enumerate}
}

\reqfield{Relevância para a aplicação}{Garante que contactos de fornecedores estão sempre à mão e que prazos de entrega são conhecidos}

\vspace{0.8cm}

% --- Requisito Stock 6 ---
\reqheader{REQ-STK-006: Estrutura de Dados de Peça}

\reqfield{Requisito de utilizador}{O sistema armazena informação completa sobre cada tipo de peça}

\reqfield{Fonte}{Documento de Análise de Requisitos}

\reqfield{Área do sistema}{Gestão de Stock}

\reqfield{Requisitos de sistema}{
	\begin{enumerate}
		\item Campo obrigatório: Fornecedor (referência a fornecedor)
		\item Campo obrigatório: Referência do fornecedor (SKU/código externo)
		\item Campo obrigatório: Preço de compra ao fornecedor (€)
		\item Campo obrigatório: Preço de venda ao cliente (€ ou incluído em reparação)
		\item Campo obrigatório: Nível mínimo aceitável em stock (quantidade)
		\item Campo opcional: Descrição da peça
		\item Campo opcional: Categoria/tipo
		\item Campo opcional: Localização física no armazém
	\end{enumerate}
}

\reqfield{Relevância para a aplicação}{Padroniza a informação de peças e permite margens corretas e encomendas baseadas em níveis mínimos}

\vspace{0.8cm}

% --- Requisito Stock 7 ---
\reqheader{REQ-STK-007: Registro de Entrada de Peças}

\reqfield{Requisito de utilizador}{O gestor de stock registar a receção de peças do fornecedor}

\reqfield{Fonte}{Documento de Análise de Requisitos}

\reqfield{Área do sistema}{Gestão de Stock}

\reqfield{Requisitos de sistema}{
	\begin{enumerate}
		\item Campo obrigatório: Referência da peça
		\item Campo obrigatório: Quantidade recebida
		\item Campo obrigatório: Fornecedor
		\item Campo obrigatório: Preço de compra unitário
		\item Campo obrigatório: Preço de venda unitário
		\item Campo obrigatório: Data de receção
		\item Campo condicional: Número de série (obrigatório se preço de compra $>$ 70€)
		\item O sistema registar automaticamente data/hora e utilizador
		\item O stock é atualizado automaticamente
	\end{enumerate}
}

\reqfield{Relevância para a aplicação}{Garante que stock é rastreável, com lote identificável para devolução ou garantia}

\vspace{0.8cm}

% --- Requisito Stock 8 ---
\reqheader{REQ-STK-008: Geração de Listas de Encomenda}

\reqfield{Requisito de utilizador}{O sistema gera automaticamente lista de peças a encomendar}

\reqfield{Fonte}{Documento de Análise de Requisitos}

\reqfield{Área do sistema}{Gestão de Stock}

\reqfield{Requisitos de sistema}{
	\begin{enumerate}
		\item O sistema calcula para cada peça: consumo médio mensal, stock atual e nível mínimo
		\item Se stock atual $<$ nível mínimo, peça é incluída na lista
		\item A quantidade a encomendar = (consumo mensal x 2) + nível mínimo - stock atual
		\item O sistema sugere fornecedor com melhor preço e prazo mais curto
		\item A lista é exportável para email/print para enviar ao fornecedor
	\end{enumerate}
}

\reqfield{Relevância para a aplicação}{Automatiza processo de encomenda, evitando stockouts e minimizando stock excessivo}

\vspace{0.8cm}

% --- Requisito Stock 9 ---
\reqheader{REQ-STK-009: Gestão de Devoluções a Fornecedor}

\reqfield{Requisito de utilizador}{O gestor de stock registar devolução de peças defeituosas ao fornecedor}

\reqfield{Fonte}{Documento de Análise de Requisitos}

\reqfield{Área do sistema}{Gestão de Stock, Devolução a Fornecedores}

\reqfield{Requisitos de sistema}{
	\begin{enumerate}
		\item O gestor pode criar registo de devolução para peças marcadas "Possível defeito"
		\item Registo inclui: peça, motivo da devolução (defeito, errado, danificado)
		\item Registo inclui: data de devolução, fornecedor contactado
		\item O estado da devolução pode ser: "Pendente devolução", "Devolvida ao fornecedor", "Inválida para devolução"
		\item Transição entre estados é registada com timestamp
		\item Consoante o estado final, stock é ajustado (aumento se devolvida, custo se inválida)
	\end{enumerate}
}

\reqfield{Relevância para a aplicação}{Recupera valor de peças defeituosas e reduz perda total do negócio}

\vspace{1.5cm}

% ======================== REQUISITOS DA ASSISTENTE ADMINISTRATIVA ========================

\section*{Requisitos da Assistente Administrativa (Secretária)}

% --- Requisito Sec 1 ---
\reqheader{REQ-SEC-001: Consulta de Ordens de Serviço por Estado}

\reqfield{Requisito de utilizador}{A secretária pode visualizar todas as OS agrupadas por estado}

\reqfield{Fonte}{Documento de Análise de Requisitos}

\reqfield{Área do sistema}{Gestão de Ordens de Serviço}

\reqfield{Requisitos de sistema}{
	\begin{enumerate}
		\item O sistema apresenta 6 listas: Pendente Diagnóstico, Pendente Reparação, A aguardar peças, Pendente aprovação, Pendente Pagamento, Paga
		\item Cada lista mostra todas as OS nesse estado
		\item A secretária pode filtrar por: estado, data de criação, nome do cliente, mecânico responsável
		\item Os resultados podem ser ordenados por data, cliente ou valor
	\end{enumerate}
}

\reqfield{Relevância para a aplicação}{Fornece visibilidade rápida do fluxo de reparações, identificando gargalos ou atrassos}

\vspace{0.8cm}

% --- Requisito Sec 2 ---
\reqheader{REQ-SEC-002: Operações CRUD de Clientes e Trotinetes}

\reqfield{Requisito de utilizador}{A secretária pode criar, editar, visualizar e remover registos de clientes e trotinetes}

\reqfield{Fonte}{Documento de Análise de Requisitos}

\reqfield{Área do sistema}{Gestão de Clientes}

\reqfield{Requisitos de sistema}{
	\begin{enumerate}
		\item A secretária pode criar novo cliente
		\item A secretária pode editar dados de cliente existente
		\item A secretária pode visualizar histórico completo de um cliente
		\item A secretária pode remover ou desativar um cliente (apenas se sem OS ativas)
		\item As mesmas operações aplicam-se a trotinetes
	\end{enumerate}
}

\reqfield{Relevância para a aplicação}{Centraliza dados de clientes no sistema, evitando dispersão e facilitando rastreamento}

\vspace{0.8cm}

% --- Requisito Sec 3 ---
\reqheader{REQ-SEC-003: Estrutura de Dados de Cliente}

\reqfield{Requisito de utilizador}{O sistema armazena informação essencial de cada cliente}

\reqfield{Fonte}{Documento de Análise de Requisitos}

\reqfield{Área do sistema}{Gestão de Clientes}

\reqfield{Requisitos de sistema}{
	\begin{enumerate}
		\item Campo obrigatório: Nome completo
		\item Campo obrigatório: Email
		\item Campo obrigatório: Telemóvel
		\item Campo obrigatório: NIF
		\item Campo opcional: Morada
		\item Campo opcional: Foto/avatar
		\item Cada cliente pode ter uma ou mais trotinetes associadas
	\end{enumerate}
}

\reqfield{Relevância para a aplicação}{Garante que dados de cliente e contacto estão sempre disponíveis para notificações e cobrança}

\vspace{0.8cm}

% --- Requisito Sec 4 ---
\reqheader{REQ-SEC-004: Estrutura de Dados de Trotinete}

\reqfield{Requisito de utilizador}{O sistema armazena especificações técnicas de cada trotinete}

\reqfield{Fonte}{Documento de Análise de Requisitos}

\reqfield{Área do sistema}{Gestão de Trotinetes}

\reqfield{Requisitos de sistema}{
	\begin{enumerate}
		\item Campo obrigatório: Marca (ex: Xiaomi, Ninebot)
		\item Campo obrigatório: Modelo (ex: M365, G30)
		\item Campo obrigatório: Número de série (único)
		\item Campo obrigatório: Tipo de motor (escova/sem escova, potência)
		\item Campo obrigatório: Cliente proprietário
		\item Campo opcional: Ano de fabrico
		\item Campo opcional: Foto/desenho
		\item Campo opcional: Histórico de problemas
	\end{enumerate}
}

\reqfield{Relevância para a aplicação}{Permite diagnóstico rápido baseado no modelo e histórico de problemas conhecidos}

\vspace{0.8cm}

% --- Requisito Sec 5 ---
\reqheader{REQ-SEC-005: Criação de Ordem de Serviço com Documentação Completa}

\reqfield{Requisito de utilizador}{A secretária pode criar uma OS com toda informação inicial necessária}

\reqfield{Fonte}{Documento de Análise de Requisitos}

\reqfield{Área do sistema}{Gestão de Ordens de Serviço}

\reqfield{Requisitos de sistema}{
	\begin{enumerate}
		\item Campo obrigatório: Cliente (seleção de cliente existente)
		\item Campo obrigatório: Trotinete (seleção de trotinete do cliente ou criação nova)
		\item Campo obrigatório: Descrição do problema relatado pelo cliente
		\item Campo opcional: Acessórios deixados (carregador, cadeado, etc.)
		\item Campo opcional: Estado ao receber (notas textuais livres)
		\item Campo opcional: Fotografias de entrada (até 5 fotos)
		\item O sistema atribui automaticamente estado "Pendente Diagnóstico"
		\item O sistema gera automaticamente referência única (ex: OS-2024-001)
		\item O sistema registar data/hora de criação e utilizador
	\end{enumerate}
}

\reqfield{Relevância para a aplicação}{Garante que diagnóstico tem contexto completo, evitando surpresas sobre estado inicial}

\vspace{0.8cm}

% --- Requisito Sec 6 ---
\reqheader{REQ-SEC-006: Consulta de Histórico de Reparações}

\reqfield{Requisito de utilizador}{A secretária pode ver rapidamente todo histórico de reparações de um cliente ou trotinete}

\reqfield{Fonte}{Documento de Análise de Requisitos}

\reqfield{Área do sistema}{Gestão de Ordens de Serviço}

\reqfield{Requisitos de sistema}{
	\begin{enumerate}
		\item A secretária pode pesquisar por nome de cliente e ver todas as suas OS
		\item A secretária pode pesquisar por número de série da trotinete
		\item O histórico mostra: data da reparação, problema, reparações feitas, peças usadas
		\item O histórico está ordenado cronologicamente (mais recentes primeiro)
		\item O histórico mostra custo e duração de cada reparação
	\end{enumerate}
}

\reqfield{Relevância para a aplicação}{Identifica padrões de problemas recorrentes e permite diagnosticar de forma mais rápida}

\vspace{0.8cm}

% --- Requisito Sec 7 ---
\reqheader{REQ-SEC-007: Registro de Aprovação de Orçamento pelo Cliente}

\reqfield{Requisito de utilizador}{A secretária pode registar quando um cliente aprova o orçamento proposto}

\reqfield{Fonte}{Documento de Análise de Requisitos}

\reqfield{Área do sistema}{Gestão de Ordens de Serviço}

\reqfield{Requisitos de sistema}{
	\begin{enumerate}
		\item O registro é possível apenas para OS em estado "Pendente aprovação do orçamento"
		\item A secretária clica botão "Aprovar orçamento"
		\item O sistema registar data/hora de aprovação
		\item O estado da OS muda automaticamente para "Pendente Reparação"
		\item Uma notificação é enviada ao mecânico responsável
	\end{enumerate}
}

\reqfield{Relevância para a aplicação}{Debloqueia a reparação e sinaliza ao mecânico que pode prosseguir}

\vspace{0.8cm}

% --- Requisito Sec 8 ---
\reqheader{REQ-SEC-008: Consulta de Estado de OS com Histórico Completo}

\reqfield{Requisito de utilizador}{A secretária pode visualizar estado atual e histórico completo de uma OS}

\reqfield{Fonte}{Documento de Análise de Requisitos}

\reqfield{Área do sistema}{Gestão de Ordens de Serviço}

\reqfield{Requisitos de sistema}{
	\begin{enumerate}
		\item Busca por nome de cliente mostra todas as OS do cliente (ordenadas por data)
		\item Busca por referência da OS mostra OS específica
		\item Busca por número de série da trotinete mostra todas as OS dessa trotinete
		\item Para cada OS, mostra estado atual
		\item Para cada OS, mostra histórico completo de mudanças de estado
		\item Histórico inclui: data/hora de cada transição, quem fez a mudança
		\item Histórico inclui: reparações realizadas, peças usadas, custos
	\end{enumerate}
}

\reqfield{Relevância para a aplicação}{Fornece informação completa para responder a questões de clientes e resolver disputas}

\vspace{0.8cm}

% --- Requisito Sec 9 ---
\reqheader{REQ-SEC-009: Registro de Data de Notificação ao Cliente}

\reqfield{Requisito de utilizador}{A secretária registar quando notificou o cliente sobre conclusão}

\reqfield{Fonte}{Documento de Análise de Requisitos}

\reqfield{Área do sistema}{Gestão de Ordens de Serviço}

\reqfield{Requisitos de sistema}{
	\begin{enumerate}
		\item Quando OS está em estado "Pendente Pagamento", a secretária registar data de notificação
		\item Esta data é essencial para cálculo de taxa de armazenagem (contagem começa daqui)
		\item O sistema apresenta um campo para entrada manual ou pode registar automaticamente
		\item Uma vez registada, a data é imutável (apenas gerente pode alterar se necessário)
	\end{enumerate}
}

\reqfield{Relevância para a aplicação}{Define o ponto de partida para acréscimo de armazenagem, sendo crítico para justiça na cobrança}

\vspace{0.8cm}

% --- Requisito Sec 10 ---
\reqheader{REQ-SEC-010: Associação de Peças em OS Pendente Pagamento}

\reqfield{Requisito de utilizador}{A secretária pode associar peças a uma OS se o mecânico não o fez}

\reqfield{Fonte}{Documento de Análise de Requisitos}

\reqfield{Área do sistema}{Gestão de Ordens de Serviço, Gestão de Stock}

\reqfield{Requisitos de sistema}{
	\begin{enumerate}
		\item A secretária pode adicionar peças apenas a OS em estado "Pendente Pagamento"
		\item A secretária procura por peça (referência, fornecedor ou descrição)
		\item Ao adicionar peça, o preço de venda é registado
		\item O custo da OS é recalculado automaticamente
		\item A quantidade em stock da peça diminui
	\end{enumerate}
}

\reqfield{Relevância para a aplicação}{Garante que se mecânico esqueceu, secretária pode registar peças antes do pagamento}

\vspace{0.8cm}

% --- Requisito Sec 11 ---
\reqheader{REQ-SEC-011: Registro de Pagamento da OS}

\reqfield{Requisito de utilizador}{A secretária registar quando cliente pagou e por qual método}

\reqfield{Fonte}{Documento de Análise de Requisitos}

\reqfield{Área do sistema}{Gestão de Ordens de Serviço, Gestão Financeira}

\reqfield{Requisitos de sistema}{
	\begin{enumerate}
		\item O registro é possível apenas para OS em estado "Pendente Pagamento"
		\item A secretária seleciona método de pagamento: numerário, multibanco, MBWay
		\item A secretária introduz data e hora do pagamento
		\item Opcionalmente, pode indicar comprovante/referência
		\item Após ambos (método selecionado + data preenchida), estado muda para "Paga"
		\item Movimento financeiro é criado automaticamente
	\end{enumerate}
}

\reqfield{Relevância para a aplicação}{Fecha o ciclo de OS e gera comprovante de conclusão}

\vspace{0.8cm}

% --- Requisito Sec 12 ---
\reqheader{REQ-SEC-012: Bloqueio de Pagamento em Cascata}

\reqfield{Requisito de utilizador}{O sistema impede que cliente pague uma OS se tem outra anterior ainda pendente}

\reqfield{Fonte}{Documento de Análise de Requisitos}

\reqfield{Área do sistema}{Gestão de Ordens de Serviço}

\reqfield{Requisitos de sistema}{
	\begin{enumerate}
		\item Antes de permitir marcação de "Paga", sistema verifica todas as OS do cliente
		\item Se cliente tem outra OS em "Pendente Pagamento" com data de notificação anterior, bloqueia
		\item Mensagem de erro indica qual OS anterior não foi paga
		\item Secretária deve resolver OS anterior antes de permitir pagamento desta
	\end{enumerate}
}

\reqfield{Relevância para a aplicação}{Evita que cliente "pule" fila de pagamento e garante ordem justa de cobrança}

\vspace{1.5cm}

% ======================== REQUISITOS DO MECÂNICO ========================

\section*{Requisitos do Mecânico}

% --- Requisito Mec 1 ---
\reqheader{REQ-MEC-001: Seleção de OS para Diagnóstico ou Reparação}

\reqfield{Requisito de utilizador}{O mecânico pode escolher qual OS deseja trabalhar entre diagnósticos ou reparações}

\reqfield{Fonte}{Documento de Análise de Requisitos}

\reqfield{Área do sistema}{Gestão de Ordens de Serviço}

\reqfield{Requisitos de sistema}{
	\begin{enumerate}
		\item A interface apresenta dois separadores/abas: "Pendente Diagnóstico" e "Pendente Reparação"
		\item Cada separador mostra lista de OS disponíveis ordenadas por data (mais antigas primeiro)
		\item O mecânico clica em uma OS para selecioná-la
		\item A OS fica bloqueada para outros mecânicos
		\item O mecânico pode desistir/devolver a OS à fila se necessário
	\end{enumerate}
}

\reqfield{Relevância para a aplicação}{Permite distribuição natural de trabalho sem necessidade de atribuição manual}

\vspace{0.8cm}

% --- Requisito Mec 2 ---
\reqheader{REQ-MEC-002: Execução de Diagnóstico com Seleção de Reparações}

\reqfield{Requisito de utilizador}{O mecânico realiza diagnóstico selecionando reparações necessárias e peças do stock}

\reqfield{Fonte}{Documento de Análise de Requisitos}

\reqfield{Área do sistema}{Gestão de Ordens de Serviço, Tabela de Reparações}

\reqfield{Requisitos de sistema}{
	\begin{enumerate}
		\item Após selecionar OS em "Pendente Diagnóstico", aparece formulário para diagnóstico
		\item O mecânico pode selecionar múltiplas reparações da tabela de reparações
		\item A tabela é consultável por referência, fornecedor ou data
		\item O mecânico pode procurar peças do stock (referência, fornecedor, descrição)
		\item O mecânico pode associar múltiplas peças à OS
		\item O mecânico pode visualizar histórico de reparações anteriores dessa trotinete
		\item O valor total do diagnóstico é calculado em tempo real
	\end{enumerate}
}

\reqfield{Relevância para a aplicação}{Estrutura o diagnóstico e documenta o que foi avaliado}

\vspace{0.8cm}

% --- Requisito Mec 3 ---
\reqheader{REQ-MEC-003: Validação de Limite de Orçamento (50€)}

\reqfield{Requisito de utilizador}{O sistema impede reparação se orçamento excede 50€ sem aprovação cliente}

\reqfield{Fonte}{Documento de Análise de Requisitos}

\reqfield{Área do sistema}{Gestão de Ordens de Serviço}

\reqfield{Requisitos de sistema}{
	\begin{enumerate}
		\item Após mecânico definir reparações e peças no diagnóstico, sistema calcula valor total
		\item Se valor total $>$ 50€, estado automático muda para "Pendente aprovação do orçamento"
		\item Se valor total $\le$ 50€, estado automático muda para "Pendente Reparação"
		\item Em caso de bloqueio, o mecânico recebe mensagem clara explicando o motivo
		\item Mecânico não pode prosseguir até aprovação cliente
	\end{enumerate}
}

\reqfield{Relevância para a aplicação}{Protege cliente de cobranças inesperadas e garante controlo de custos}

\vspace{0.8cm}

% --- Requisito Mec 4 ---
\reqheader{REQ-MEC-004: Execução de Reparação com Seleção de Serviços Realizados}

\reqfield{Requisito de utilizador}{O mecânico registar quais reparações foram efetivamente executadas}

\reqfield{Fonte}{Documento de Análise de Requisitos}

\reqfield{Área do sistema}{Gestão de Ordens de Serviço}

\reqfield{Requisitos de sistema}{
	\begin{enumerate}
		\item Após selecionar OS em "Pendente Reparação", aparece formulário de execução
		\item O mecânico pode selecionar múltiplas reparações da tabela (as que efetivamente fez)
		\item O mecânico pode procurar e associar peças usadas (referência, fornecedor, descrição)
		\item O mecânico pode visualizar histórico de reparações anteriores
		\item O valor final é recalculado baseado nas reparações e peças reais
		\item Pode haver diferenças entre diagnóstico previsto e execução real
	\end{enumerate}
}

\reqfield{Relevância para a aplicação}{Registar efetivamente o que foi feito, não apenas o que foi diagnosticado}

\vspace{0.8cm}

% --- Requisito Mec 5 ---
\reqheader{REQ-MEC-005: Requisicao de Pecas em Falta com Estado Aguardando Pecas}

\reqfield{Requisito de utilizador}{O mecânico pode requisitar peças em falta sem suspender a reparação}

\reqfield{Fonte}{Documento de Análise de Requisitos}

\reqfield{Área do sistema}{Gestão de Ordens de Serviço, Gestão de Stock}

\reqfield{Requisitos de sistema}{
	\begin{enumerate}
		\item Enquanto mecânico está a executar OS "Pendente Reparação", pode adicionar peça em falta
		\item Sistema detecta que peça não tem stock disponível
		\item O mecânico pode clicar "Requisitar encomenda" para essa peça
		\item A OS automaticamente muda para estado "A aguardar peças"
		\item O gestor de stock recebe alerta com peça requerida
		\item A reparação pode continuar normalmente com outras peças/serviços
		\item Quando peça chega e é adicionada ao stock, OS volta a "Pendente Reparação"
	\end{enumerate}
}

\reqfield{Relevância para a aplicação}{Evita interrupção completa da reparação por falta de uma peça, acelerando conclusão}

\vspace{0.8cm}

% --- Requisito Mec 6 ---
\reqheader{REQ-MEC-006: Identificação de Novos Problemas Durante Reparação}

\reqfield{Requisito de utilizador}{O mecânico pode adicionar novos problemas descobertos após diagnóstico inicial}

\reqfield{Fonte}{Documento de Análise de Requisitos}

\reqfield{Área do sistema}{Gestão de Ordens de Serviço}

\reqfield{Requisitos de sistema}{
	\begin{enumerate}
		\item Durante execução, o mecânico pode clicar "Adicionar novo problema"
		\item Seleciona problema/reparação da tabela
		\item Se custo total exceder 50€, OS muda automaticamente para "Pendente aprovação do orçamento"
		\item Uma notificação é enviada à secretária com aviso de custo adicional
		\item Mecânico não pode prosseguir até novo orçamento ser aprovado
	\end{enumerate}
}

\reqfield{Relevância para a aplicação}{Recompõe orçamento se problemas adicionais aparecem, mantendo cliente informado}

\vspace{0.8cm}

% --- Requisito Mec 7 ---
\reqheader{REQ-MEC-007: Reportar Defeito em Peça Instalada}

\reqfield{Requisito de utilizador}{O mecânico pode reportar que uma peça instalada é defeituosa}

\reqfield{Fonte}{Documento de Análise de Requisitos}

\reqfield{Área do sistema}{Gestão de Ordens de Serviço, Gestão de Stock}

\reqfield{Requisitos de sistema}{
	\begin{enumerate}
		\item Após adicionar peça à lista, o mecânico pode marcar como "Possível defeito"
		\item A peça é removida da lista de peças usadas
		\item O custo da OS diminui (preço da peça é subtraído)
		\item A peça é marcada para devolução ao fornecedor
		\item Um alerta é gerado para o gestor de stock
	\end{enumerate}
}

\reqfield{Relevância para a aplicação}{Evita cobrar cliente por peça defeituosa e dispara investigação de fornecedor}

\vspace{0.8cm}

% --- Requisito Mec 8 ---
\reqheader{REQ-MEC-008: Checklist de Segurança Obrigatória}

\reqfield{Requisito de utilizador}{O mecânico deve completar checklist de segurança antes de concluir reparação}

\reqfield{Fonte}{Documento de Análise de Requisitos}

\reqfield{Área do sistema}{Gestão de Ordens de Serviço}

\reqfield{Requisitos de sistema}{
	\begin{enumerate}
		\item Antes de marcar OS como concluída, um formulário de checklist aparece
		\item Checklist inclui: verificação de travões, luzes, pneus, aceleração, travagem, visor
		\item Checklist inclui: teste prático de condução
		\item O mecânico marca cada item como OK/falho
		\item Se algum item falha, reparação não pode ser concluída (precisa de correção)
		\item Todos os itens devem estar OK para prosseguir
	\end{enumerate}
}

\reqfield{Relevância para a aplicação}{Garante que trotinete está segura para uso antes de devolver ao cliente, reduzindo acidentes}

\vspace{0.8cm}

% --- Requisito Mec 9 ---
\reqheader{REQ-MEC-009: Marcação de Conclusão e Transição para Pagamento}

\reqfield{Requisito de utilizador}{O mecânico marca a OS como concluída, iniciando fluxo de pagamento}

\reqfield{Fonte}{Documento de Análise de Requisitos}

\reqfield{Área do sistema}{Gestão de Ordens de Serviço}

\reqfield{Requisitos de sistema}{
	\begin{enumerate}
		\item Após passar checklist, o mecânico clica "Concluir reparação"
		\item O estado da OS muda automaticamente para "Pendente Pagamento"
		\item A OS é desbloqueada (outro mecânico não está mais bloqueado por ela)
		\item Um alerta é gerado para a secretária comunicar conclusão ao cliente
		\item O sistema registar data/hora de conclusão
	\end{enumerate}
}

\reqfield{Relevância para a aplicação}{Marca fim do trabalho operacional, iniciando fase de cobrança}

\vspace{1.5cm}

% ======================== REQUISITOS DE INVESTIGAÇÃO OPERACIONAL ========================

\section*{Requisitos de Investigação Operacional}

% --- Requisito IO 1 ---
\reqheader{REQ-IO-001: Registro de Danos Prévios com Fotografias}

\reqfield{Requisito de utilizador}{O sistema permite registar rapidamente danos pré-existentes na trotinete ao receber}

\reqfield{Fonte}{Documento de Análise de Requisitos}

\reqfield{Área do sistema}{Gestão de Ordens de Serviço}

\reqfield{Requisitos de sistema}{
	\begin{enumerate}
		\item Ao criar OS, há secção dedicada "Registo de Danos Prévios"
		\item Esta secção permite descrição textual de danos visíveis
		\item Esta secção permite anexação de múltiplas fotografias (até 5)
		\item As fotos são armazenadas na ficha da OS para referência futura
		\item Estabelece baseline de condição inicial da trotinete
	\end{enumerate}
}

\reqfield{Relevância para a aplicação}{Protege a oficina contra reclamações falsas de dano pós-reparação}

\vspace{0.8cm}

% --- Requisito IO 2 ---
\reqheader{REQ-IO-002: Registro de Acessórios Deixados pelo Cliente}

\reqfield{Requisito de utilizador}{O sistema permite nomear e registar acessórios deixados pelo cliente}

\reqfield{Fonte}{Documento de Análise de Requisitos}

\reqfield{Área do sistema}{Gestão de Ordens de Serviço}

\reqfield{Requisitos de sistema}{
	\begin{enumerate}
		\item Ao criar OS, há campo para "Itens deixados pelo cliente"
		\item Exemplos pré-definidos: carregador, cadeado, chave, mochila, outros
		\item O sistema permite checkboxes para itens comuns
		\item O sistema permite campo livre para outros itens
		\item Ao devolver trotinete, sistema verifica que todos os itens foram devolvidos
		\item Se falta algum item, alertar antes de libertar trotinete
	\end{enumerate}
}

\reqfield{Relevância para a aplicação}{Evita disputas sobre itens perdidos e responsabilização incorreta}

\vspace{0.8cm}

% --- Requisito IO 3 ---
\reqheader{REQ-IO-003: Quadro de Estado em Tempo Real na Receção}

\reqfield{Requisito de utilizador}{Na receção há um quadro visível com estado atual de todas as OS}

\reqfield{Fonte}{Documento de Análise de Requisitos}

\reqfield{Área do sistema}{Interface de Utilizador, Gestão de Ordens de Serviço}

\reqfield{Requisitos de sistema}{
	\begin{enumerate}
		\item Na receção há um ecrã ou tablet com interface dedicada
		\item O mecânico introduz referência da OS ou scans QR/barcode
		\item Sistema mostra estado atual da OS em grande
		\item O mecânico pode confirmar "ainda em diagnóstico", "em reparação", "à espera de peças", "pronta para levantar"
		\item Estado é atualizado imediatamente no sistema
		\item A secretária vê atualizações em tempo real
		\item Dispensa necessidade de chamadas/comunicação manual entre oficina e receção
	\end{enumerate}
}

\vspace{0.8cm}

\end{document}
